\section{Related Work}

\subsection{Medical vision-language models for chest X-ray}

Medical VLMs have rapidly expanded from task-specific classifiers to instruction-following multimodal models. CheXagent introduced a CXR-specific foundation model trained on CheXinstruct, a large multi-task CXR corpus, and demonstrated strong performance across perception, reasoning, and report-generation tasks \citep{chen2024chexagent}. MedGemma extends the Gemma family with medical multimodal pretraining and instruction tuning, including chest X-ray capabilities \citep{sellergren2025medgemma}. LLaVA-Med brought instruction tuning to biomedical vision-language tasks on a small compute budget \citep{li2024llavamed}. General-domain models such as Qwen2.5-VL and InternVL provide strong open-weight visual recognition and structured output capabilities without medical-specific fine-tuning \citep{bai2025qwen25vl}. At the same time, negative results have shown that vision-language pretraining does not necessarily produce transferable representations for dense multi-label CXR classification \citep{fisher2025negative}. This tension between broad multimodal capability and fine-grained medical reliability motivates our joint evaluation of discrimination, calibration, and image reliance.

\subsection{Calibration and uncertainty in medical AI}

Calibration measures whether predictive confidence matches empirical correctness. In medical imaging, miscalibration can be clinically dangerous because overconfident errors may escape human review. CALM-VLM integrated temperature scaling and selective abstention into a generative 3D VLM for brain MRI classification, improving coverage-adjusted ROC-AUC across multi-site data \citep{dhinagar2026calm}. Calibration for medical VQA has also been studied empirically and through verbalised uncertainty training (MICCAI 2025; Just how sure are you, 2026). However, these works primarily target single-modality or internal settings and do not test whether calibration transfer survives image-quality shift. Our work evaluates calibration after source-domain fitting and applies it to shifted and external CXR data, while explicitly conditioning reliability on visual evidence use.

\subsection{Selective prediction and abstention}

Selective prediction allows a model to abstain on uncertain cases, trading coverage for lower risk. Recent work has formalised fine-grained abstention for medical VLMs, including answer, clarify, and abstain actions under image corruption, text corruption, and cross-modal mismatch \citep{nguyen2026answer}. AT-CXR combined uncertainty estimates with an agentic triage policy on NIH ChestX-ray14 and achieved strong risk--coverage performance for pulmonary edema triage \citep{li2025atcxr}. These results establish the promise of selective prediction. Our work differs in two ways. First, we evaluate multi-label CXR findings under a systematic image-quality shift benchmark rather than a single balanced triage task. Second, we show that selective risk--coverage can appear favourable for a model that is image-invariant, motivating image-reliance-aware selective trust.

\subsection{Distribution shift and robustness}

Distribution shift in medical imaging arises from differences in scanners, acquisition protocols, patient populations, and post-processing. Prior work has documented domain shift and dataset bias in public CXR datasets \citep{rafferty2026limitations}, and robustness studies have shown that VLMs are sensitive to medical image artefacts \citep{cheng2025robustness}. Stress-testing studies have further revealed that multimodal medical benchmarks can overstate readiness when critical visual inputs are removed \citep{gu2025illusion}. Our shift benchmark complements these efforts by measuring performance, calibration, and image sensitivity across six image-quality axes and three severity levels. A further concern we raise is that robustness is usually reported as aggregate accuracy, which on a label-imbalanced task can hide degradation concentrated in the minority pathology labels; we report positive recall alongside accuracy for this reason.

\subsection{Image reliance, grounding, and hallucination}

A central question for medical VLMs is whether predictions are grounded in image content or driven by language and label priors. Studies using blank-image substitution have found substantial variability in visual dependency across frontier models \citep{felizzi2025grounded}. Large-scale analyses of multimodal medical models have shown that image predictions can be dominated by text and that misleading clinical vignettes can sharply reduce accuracy \citep{buckley2026exploit}. Counterfactual evaluations with real, blank, and shuffled images have introduced metrics such as visual reliance score, image sensitivity, and hallucinated visual reasoning rate \citep{zafar2026beyond}. For radiology report generation, hallucination benchmarks and mitigation methods have begun to characterise factual inconsistencies (Heiman et al., 2025; MedVH, 2025). Our work connects this literature to selective trust: we treat image reliance as a first-class reliability signal, measure it as a continuous spectrum across seven models, and show that an image-blind model can simultaneously produce valid structured outputs and apparently acceptable selective metrics.

\subsection{Causal image-use audits and the detection gap}

A second, more recent line of work tests image use by intervention rather than by benchmark accuracy. Lotfinia et al. (2026) audit nine CXR systems with three interventions --- swapping in a same-label patient, masking the radiologist-marked target region, and masking an irrelevant region --- and read three jointly defined quantities (a causal grounding rate, an unrelated-image answer rate, and an irrelevant-mask stability) to sort models into image-using, image-ignoring, and unstable regimes; a text-only model lands within 5.7 accuracy points of the best multimodal system, and image-ignoring models report high confidence on their errors \citep{lotfinia2026image}. MedVIGIL frames the same concern as an evidence contract, operationalising silent failure under perturbed evidence as the central failure mode in a 300-case, four-radiologist suite over text- and image-side probes \citep{jiang2026medvigil}.

For selective trust, the consequential question is whether such grounding failures can be detected automatically. Khanmohammadi et al. (2026) evaluate nine confidence estimators --- including a grounding-aware estimator --- and report that none reliably detects when the image was ignored, making this shortfall a central negative finding \citep{khanmohammadi2026calibrated}. By contrast, Chowdhury and Miah (2026) show that a model-agnostic external probe on a frozen CLIP vision encoder, measuring text-conditioned layer-wise patch alignment, reduces the mirage rate on blank, unrelated, and mismatched inputs from 21.7--66.6\% down to 2.7--3.3\% across twelve backbones \citep{chowdhury2026mirage}. The gap between these two results --- internal estimators that fail and an external encoder that succeeds --- leaves open whether a model's own probability alignment under counterfactual inputs can detect image ignorance without a separate vision encoder. Our visual-reliance signal addresses exactly this question, and we report its error-versus-coverage trade-off rather than a single accuracy figure.

\subsection{Positioning}

Existing work has studied medical VLM calibration, selective prediction, robustness, grounding, and causal image use largely in isolation. Calibration and selective prediction have been advanced on MRI and CXR triage (Dhinagar et al., 2026; Li et al., 2025; Nguyen et al., 2026); robustness and grounding studies have shown that image quality and visual dependency matter (Cheng et al., 2025; Buckley et al., 2026; Zafar et al., 2026); and causal audits have established that image use is graded and that internal estimators cannot detect its absence (Lotfinia et al., 2026; Khanmohammadi et al., 2026). We contribute by jointly evaluating calibration, uncertainty, selective prediction, shift response, positive recall, image reliance, and failure mode across open-source CXR VLMs, at the granularity of individual labels rather than whole answers, and across three layers --- synthetic in-domain shift, cross-hospital native shift, and free-text report generation. The resulting failures form a coherent spectrum rather than isolated phenomena. This motivates the shift-aware selective trust framework developed in this paper.